%% 11_future_directions.tex - Future Directions

\section{Future Directions}
\label{sec:future-directions}

The field of LLM evaluation is rapidly evolving. This section highlights emerging trends and open challenges that will shape evaluation practices in the coming years.

\subsection{Standardization of Application-Level Benchmarks}

Most public benchmarks, including MMLU, HELM, and BIG-Bench, evaluate foundation model capabilities rather than the performance of deployed applications~\citep{hendrycks2021measuring, liang2023holistic, srivastava2023beyond}. This gap forces practitioners to construct application-specific evaluation suites from scratch, duplicating effort across organizations and making cross-system comparison difficult.

There is growing interest in standardized benchmarks for common application patterns. Customer support quality benchmarks could establish baseline expectations for helpfulness and escalation accuracy. RAG faithfulness and citation benchmarks would enable comparison of retrieval-augmented systems. Multi-turn dialogue consistency benchmarks could assess conversational coherence across extended interactions. Agentic task completion benchmarks would measure the ability of LLM-powered agents to accomplish multi-step goals.

Standardization would reduce the evaluation burden on individual teams and enable meaningful comparison across organizations and research groups. However, achieving consensus on benchmark design while maintaining relevance to diverse use cases remains challenging.

\subsection{Evaluation Harnesses and Frameworks}

Open-source evaluation frameworks have matured significantly. The lm-evaluation-harness~\citep{eleutherai2023lmeval} and OpenAI Evals~\citep{openai2023evals} provide infrastructure for running automated evaluations at scale. Specialized frameworks like RAGAS~\citep{es2024ragas} and ARES~\citep{saadfalcon2024ares} target retrieval-augmented generation specifically.

Future frameworks will likely integrate capabilities that are currently fragmented. A critical open challenge is meta-evaluation: benchmarking the metrics themselves. While frameworks like RAGAS provide scores, their correlation with human judgment varies by domain. Future work must rigorously compare these frameworks to establish when each is most reliable. Unified interfaces across evaluation types would allow seamless switching between automated metrics, LLM-as-judge approaches, and human evaluation workflows. Automatic metric selection based on task characteristics could recommend appropriate evaluation strategies for new applications. Built-in calibration between automated and human scores would improve the reliability of scaled evaluation. Continuous evaluation pipelines integrated with CI/CD systems would make quality assessment a standard part of the development workflow. Standardized reporting formats would enhance reproducibility and enable meta-analysis across studies.

\subsection{Observability and Production Monitoring}

LLM observability is nascent compared to traditional software monitoring. Few organizations have robust pipelines for assessing output quality in production beyond basic error rates and latency metrics.

Purpose-built observability platforms are emerging with capabilities specifically designed for LLM applications. Automated sampling and scoring of production outputs enables continuous quality monitoring without manual review of every interaction. Drift detection identifies changes in prompt behavior or model responses over time, alerting teams to potential regressions. Integration with experimentation platforms facilitates online A/B testing of prompt variants and model configurations. Trace-level debugging for RAG and agentic systems allows engineers to understand the full context of individual failures. Real-time alerting for quality degradation enables rapid response to emerging issues.

Investment in observability infrastructure will become increasingly critical as LLM applications proliferate and quality expectations rise.

\subsection{Agentic and Multi-Step Evaluation}

Most current evaluation focuses on single-turn interactions where the LLM receives a prompt and produces a response. Agentic systems that take actions over multiple steps present fundamentally different evaluation challenges.

New evaluation paradigms are emerging to address these challenges. Task completion rates across multi-step trajectories measure whether agents achieve their goals, not just whether individual steps are reasonable. Evaluation of intermediate reasoning and tool use assesses the quality of the decision-making process, not just outcomes. Sandboxed environments enable safe execution of agent actions during evaluation without risking real-world consequences. Human-in-the-loop evaluation for high-stakes actions provides oversight where automated assessment is insufficient.

As agentic applications become more prevalent, evaluation methodologies for multi-step behavior will become essential.

\subsection{Automated Red-Teaming}

Red-teaming is currently a largely manual process, relying on human creativity to discover failure modes and adversarial inputs~\citep{perez2022red, ganguli2022red}. This approach is expensive and does not scale well.

Automated red-teaming approaches are gaining traction. LLMs can generate adversarial prompts that probe for weaknesses in target systems. Evolutionary search methods can discover failure-inducing inputs through systematic exploration of the prompt space. Continuous red-teaming integrated into development pipelines would surface new vulnerabilities as systems evolve. Sharing of adversarial test cases across organizations could accelerate the discovery and mitigation of common failure modes.

\subsection{Better LLM Judges}

Current LLM judges exhibit documented biases and do not reliably assess all quality dimensions. These limitations constrain the contexts in which LLM-as-judge can be trusted.

Several directions show promise for improvement. Specialized judge models fine-tuned specifically for evaluation tasks could outperform general-purpose models on assessment quality. Ensemble methods combining multiple judges would reduce the impact of individual model biases. Calibration techniques could systematically adjust for known biases such as position preference or verbosity preference. Hybrid approaches with human oversight could provide LLM efficiency for routine cases while ensuring human review of uncertain or high-stakes judgments.

\subsection{Personalization and Context-Dependent Evaluation}

As LLM applications become more personalized, evaluation must account for user-specific factors. Quality definitions may vary based on user preferences, expertise levels, and interaction history. Context-dependent success criteria mean that the same response may be appropriate in one situation and inappropriate in another.

Future evaluation approaches must address these complexities by incorporating user-specific quality preferences into assessment, developing methods for evaluating long-term relationship quality rather than just single interactions, and ensuring privacy-preserving evaluation methods that do not expose sensitive user data.

